% ----------------------------------------------------------------------
%  1.1. Кристаллическое состояние
% ----------------------------------------------------------------------
\begin{orientir}
Что такое кристалл с точки зрения внешнего вида (законы огранки), с
точки зрения свойств (однородность, анизотропия, симметрия) и с точки
зрения внутреннего строения (решётка и структура). Главный итог:
\vazhno{структура = решётка + базис}.
\end{orientir}

\begin{opredelenie}[Кристалл]
Твёрдое тело с упорядоченной трёхмерно-периодической атомной
структурой. Кристаллическое состояние~— термодинамически равновесное
состояние твёрдого тела: каждой твёрдой фазе данного состава при данных
внешних условиях отвечает одна определённая кристаллическая структура.
\end{opredelenie}

Естественные плоские грани кристалла (огранка)~— внешнее проявление его
внутреннего атомного строения. Именно изучая огранку, кристаллографы
пришли к идее периодичности задолго до того, как удалось «увидеть» атомы.

\subsection{Основные законы кристаллографии}\label{ssec:laws}

Законы кристаллографии получены из наблюдений за огранкой, поэтому их
называют геометрическими (эмпирическими). Внешний облик кристаллов~—
\term{габитус}: размеры, форма и даже число граней~— у разных кристаллов
одного вещества может быть разным. Но их всегда можно расположить так,
что некоторые грани окажутся параллельными; такие грани называют
\term{соответственными}.

\begin{glavnoe}[Закон постоянства углов (Н.~Стенон, 1669)]
Углы между соответственными гранями кристаллов данного вещества
постоянны и характерны для этого вещества.
\end{glavnoe}

Стенон установил закон на кварце и гематите. Форма граней может сильно
меняться в зависимости от условий роста (рис.~\ref{fig:steno}), но углы
остаются прежними~— поэтому по углам, измеренным оптическим
\emph{гониометром}, определяют вещество.

\begin{figure}[ht]
\centering
\begin{tikzpicture}[scale=0.95]
  \def\hexR{1.2}
  \draw[fmBlue, thick, fill=fmBlue!8] (0:\hexR) \foreach \a in {60,120,...,300}{-- (\a:\hexR)} -- cycle;
  \node[font=\small] at (0,-1.55) {\textit{а}};
  \begin{scope}[xshift=4.3cm]
    \coordinate (P1) at (0:1.9); \coordinate (P2) at ($(P1)+(120:0.7)$);
    \coordinate (P3) at ($(P2)+(180:2.4)$); \coordinate (P4) at ($(P3)+(240:1.6)$);
    \coordinate (P5) at ($(P4)+(300:0.7)$); \coordinate (P6) at ($(P5)+(0:2.4)$);
    \draw[fmBlue, thick, fill=fmBlue!8] (P1) -- (P2) -- (P3) -- (P4) -- (P5) -- (P6) -- cycle;
    \foreach \p in {P1,P2,P3,P4,P5,P6}{\fill[fmRed] (\p) circle (1.2pt);}
    \node[font=\scriptsize, fmRed] at ($(P2)+(-0.15,-0.35)$) {$120^\circ$};
    \node[font=\scriptsize, fmRed] at ($(P5)+(0.15,0.35)$) {$120^\circ$};
    \node[font=\small] at (0.3,-1.55) {\textit{б}};
  \end{scope}
\end{tikzpicture}
\caption{Закон постоянства углов: поперечные сечения двух кристаллов
кварца~— «идеального» (\textit{а}) и искажённого при росте (\textit{б}).
Длины граней разные, но углы между соответственными гранями одинаковы
(здесь все по $120^\circ$)}
\label{fig:steno}
\end{figure}

\begin{glavnoe}[Закон рациональных параметров (Р.\,Ж.~Гаюи, 1784)]
Если за оси координат принять три ребра кристалла, то отрезки, которые
любая грань отсекает на осях, относятся как \emph{целые} числа, если их
выразить в некоторых осевых единицах $a$, $b$, $c$ (свои для каждого
вещества).
\end{glavnoe}

Существование осевых единиц, одинаковых для всех граней, означает, что
кристалл построен из одинаковых «кирпичиков»: вместе с представлением об
атомах закон Гаюи с необходимостью ведёт к выводу о \vazhno{трёхмерной
периодичности атомного строения}.

Дальнейшие вехи:
\begin{itemize}
\item Е.\,С.~Фёдоров (1890) и А.~Шёнфлис (1891) теоретически вывели все
\emph{230 пространственных групп} симметрии периодических структур
(п.~1.3.6);
\item первым прямым доказательством существования пространственной
решётки стала \emph{дифракция рентгеновских лучей} на кристаллах
(М.~Лауэ, В.~Фридрих, П.~Книппинг, 1912);
\item в середине XX~в. понятие симметрии обобщили: в группах
\emph{антисимметрии} (А.\,В.~Шубников) и \emph{цветной симметрии}
(Н.\,В.~Белов) кроме трёх пространственных координат есть четвёртая
«координата» иной физической природы~— время, направление спина, знак
заряда. Такие группы (шубниковские и беловские) описывают, например,
магнитные структуры.
\end{itemize}

\subsection{Макроскопические характеристики кристаллов}\label{ssec:macro}

С макроскопической точки зрения кристалл обладает тремя главными
особенностями: однородностью, анизотропией и симметрией. Пусть $F$~—
какое-то свойство: скалярное (плотность, теплоёмкость), векторное
(поляризация) или в общем случае тензорное (упругость).

\paragraph{Однородность}
Свойства во всех участках кристалла одинаковы:
\begin{equation}\label{eq:homog}
F(\vb x)=F(\vb x+\vb x') ,
\end{equation}
т.\,е. свойства не меняются при произвольном переносе начала координат.
Это позволяет рассматривать кристалл как \emph{сплошную среду} и описывать
многие свойства феноменологически, не вспоминая об атомах. Однородны,
впрочем, и жидкости, и газы, и аморфные тела.

\paragraph{Анизотропия}
Отличает кристалл от этих сред зависимость свойств от направления $\vb n$,
вдоль которого их измеряют:
\begin{equation}\label{eq:anis}
F(\vb n_1)\ne F(\vb n_2) .
\end{equation}
Анизотропны также текстуры, жидкие кристаллы, полимеры~— анизотропия
всегда определяется строением вещества.

\begin{primer}[Анизотропия упругости железа]
Модуль Юнга монокристалла железа вдоль ребра куба
$E_{\nfam{100}}\approx131$~ГПа, а вдоль пространственной диагонали
$E_{\nfam{111}}\approx283$~ГПа~— более чем вдвое больше
(рис.~\ref{fig:young-polar}). Для сравнения: вольфрам упруго почти
изотропен ($E\approx410$~ГПа во всех направлениях), хотя тоже имеет
кубическую решётку.
\end{primer}

\begin{figure}[ht]
\centering
\begin{tikzpicture}
\begin{axis}[width=0.62\textwidth, height=0.62\textwidth, axis equal image,
  xmin=-460, xmax=460, ymin=-460, ymax=460, axis lines=none,
  legend style={at={(1.0,0.0)}, anchor=south west, font=\scriptsize}]
  \foreach \r in {100,200,300,400}{\edef\tmp{\noexpand\draw[fmGray!40, thin] (axis cs:0,0) circle[radius=\r];}\tmp}
  \draw[fmGray!60] (axis cs:-440,0) -- (axis cs:440,0) (axis cs:0,-440) -- (axis cs:0,440);
  \draw[fmGray!60, dashed] (axis cs:-380,-268.7) -- (axis cs:380,268.7);
  \draw[fmGray!60, dashed] (axis cs:380,-268.7) -- (axis cs:-380,268.7);
  \addplot[fmRed, thick] table[x=xf, y=yf] {figures/cr-young-polar.dat};
  \addlegendentry{Fe}
  \addplot[fmBlue, thick, dashed] table[x=xw, y=yw] {figures/cr-young-polar.dat};
  \addlegendentry{W}
  \node[font=\scriptsize, anchor=south] at (axis cs:0,440) {$[001]$};
  \node[font=\scriptsize, anchor=west] at (axis cs:440,0) {$[110]$};
  \node[font=\scriptsize, anchor=south west] at (axis cs:380,268.7) {$[111]$};
  \foreach \r in {100,200,300,400}{\edef\tmp{\noexpand\node[font=\noexpand\tiny, fmGray, anchor=north west] at (axis cs:\r,0) {\r};}\tmp}
\end{axis}
\end{tikzpicture}
\caption{Полярная диаграмма модуля Юнга (ГПа) монокристаллов железа и
вольфрама в плоскости $(1\br10)$: длина радиус-вектора равна модулю в
данном направлении (расчёт по упругим постоянным, п.~1.5.4)}
\label{fig:young-polar}
\end{figure}

\paragraph{Симметрия}
Анизотропия не означает, что свойства различны во \emph{всех}
направлениях. Наоборот, для некоторых направлений они закономерно
одинаковы (на рис.~\ref{fig:young-polar} модуль одинаков для всех
направлений типа $\nfam{111}$). Это и есть симметрия:
\begin{equation}\label{eq:symm}
F(\vb x)=F(\vb x')=\ldots=F(\vb x^{(n)}) ,
\end{equation}
т.\,е. функция $F$ инвариантна относительно определённых преобразований
аргумента. На языке симметрии однородность~— инвариантность
относительно любого параллельного переноса, а анизотропия проявляется
\emph{в рамках} симметрии: сами функции, описывающие свойства,
симметричны.

\begin{glavnoe}
По макроскопическим признакам кристаллическое вещество~— это
\vazhno{однородная анизотропная симметричная среда}.
\end{glavnoe}

\subsection{Решётка и структура кристаллов}\label{ssec:lattice-structure}

Конкретное расположение атомов в кристалле называют его
\term{структурой}. Любой кристалл можно разбить на одинаковые ячейки,
содержащие один или несколько определённым образом расположенных атомов;
бесконечное периодическое повторение таких ячеек~— \term{кристаллическое
пространство}.

\begin{opredelenie}[Преобразование симметрии и гомологичные точки]
\term{Преобразование симметрии}~— операция, переводящая точки
кристаллического пространства в новые положения без изменения расстояний
между ними, после которой пространство нельзя отличить от исходного ни
по свойствам, ни по геометрии. Две точки $O$ и $O'$ \term{гомологичны}
(эквивалентны), если существует преобразование симметрии, переводящее $O$
в $O'$.
\end{opredelenie}

Математически кристаллическое пространство строится на трёх
постулатах:
\begin{enumerate}
\item не все точки пространства одинаковы;
\item пространство \emph{однородно}: существует шар постоянного радиуса $R$
(шар однородности), внутри которого, где бы его ни поместить, найдётся
точка, гомологичная любой заданной; рентгеновские данные показывают, что
$R$ не превышает нескольких нанометров;
\item пространство \emph{дискретно}: есть точки, вокруг которых в шаре
некоторого радиуса $r$ (шар дискретности) нет ни одной гомологичной им
точки.
\end{enumerate}
Этих постулатов достаточно, чтобы построить всю геометрическую
кристаллографию.

\paragraph{Узловой ряд и решётка}
Возьмём любую точку $O$ и шар однородности вокруг неё. В шаре найдётся
хотя бы одна точка, гомологичная $O$; ближайшую из них обозначим $1$, а
расстояние до неё~— $a_1$. Раз точка $1$ ничем не отличается от $O$, а у
$O$ справа на расстоянии $a_1$ есть сосед $1$, то и у точки $1$ на том же
расстоянии справа есть гомологичная точка $2$, и так до бесконечности; то
же влево. Получается бесконечный ряд гомологичных точек~— \term{узловой
ряд} с периодом $a_1$, называемым \term{длиной трансляции}. Значит,
\emph{трансляция}~— перенос на $a_1$~— есть преобразование симметрии
бесконечного кристаллического пространства. Повторив рассуждение в двух
других направлениях, получим, что гомологичные точки образуют
\term{пространственную решётку}.

\begin{figure}[ht]
\centering
\begin{tikzpicture}[x=1cm, y=1cm]
  \foreach \i in {0,1,2,3}{
    \node[atom=4.2mm] at (2.6*\i,0) {\scriptsize A};
    \node[atom=3.2mm, fill=atomB] at (2.6*\i+0.85,0) {\scriptsize B};
    \node[atom=2.6mm, fill=atomI] at (2.6*\i+1.7,0) {\tiny C};}
  \node[font=\small, anchor=east] at (-0.5,0) {\textit{а}};
  \foreach \i in {0,1,2,3}{\fill[fmBlue] (2.6*\i,-1.1) circle (2pt);}
  \draw[fmBlue] (-0.3,-1.1) -- (8.4,-1.1);
  \draw[dimline] (0,-1.45) -- node[below, font=\scriptsize]{$a_1$} (2.6,-1.45);
  \node[font=\small, anchor=east] at (-0.5,-1.1) {\textit{б}};
  \foreach \i in {0,1,2,3}{\draw[fmRed] (2.6*\i+0.425,-2.3) node[cross out, draw, inner sep=1.6pt]{};}
  \draw[fmRed] (-0.3,-2.3) -- (8.4,-2.3);
  \node[font=\small, anchor=east] at (-0.5,-2.3) {\textit{в}};
\end{tikzpicture}
\caption{Структура и решётка: \textit{а}~— одномерная структура из атомов
трёх сортов; \textit{б}~— узловой ряд с узлами в центрах атомов A;
\textit{в}~— тот же ряд, сдвинутый так, что узлы оказались посередине
между атомами A и B. Положение узлов произвольно, период один и тот же}
\label{fig:row}
\end{figure}

Разницу между структурой и решёткой хорошо видно на одномерном примере
(рис.~\ref{fig:row}). Поместив узел в центр атома A, найдём остальные узлы
в центрах всех атомов A. Поместив его посередине между A и B, найдём
остальные узлы в таких же серединах. Узлы~— не атомы, а гомологичные
точки, и положение начального узла выбирается произвольно.

\begin{glavnoe}[Структура = решётка + базис]
Структура кристалла~— физическая реальность~— состоит из атомов.
Пространственная решётка~— математическая абстракция, способ описать
\emph{периодичность}: она состоит из узлов. Чтобы из решётки получить
структуру, в каждый узел надо «повесить» одинаковую группу атомов~—
\term{базис} (на рис.~\ref{fig:row} базис~— тройка атомов A, B, C).
\end{glavnoe}

\begin{itogi}
\begin{itemize}
\item Кристалл~— твёрдое тело с трёхмерно-периодической структурой;
кристаллическое состояние равновесно.
\item Законы постоянства углов (Стенон) и рациональных параметров (Гаюи)
привели к идее периодичности; её прямое доказательство~— дифракция
рентгеновских лучей (Лауэ, 1912); всего 230 пространственных групп
(Фёдоров, Шёнфлис).
\item Макроскопически кристалл~— однородная анизотропная симметричная
среда.
\item Гомологичные точки образуют пространственную решётку; структура =
решётка + базис.
\end{itemize}
\end{itogi}

\begin{proverka}
\item Почему кристаллы одного вещества могут иметь разную форму, но
одинаковые углы между гранями?
\item Какой вывод о строении кристаллов следует из закона рациональных
параметров?
\item Чем однородность отличается от изотропии? Приведите пример
однородной, но анизотропной среды.
\item Почему узлы решётки не обязательно совпадают с атомами?
\item Сколько атомов в базисе структуры на рис.~\ref{fig:row}?
\end{proverka}
